\label{sec:results}

Table~\ref{tab:incumbency-summary} presents the three incumbency metrics for
\textsc{bisect} algorithmic maps, enacted 2022 congressional maps, and the
incumbency-neutral random baseline for NC, WI, and TX.

\begin{table}[ht]
\centering
\begin{threeparttable}
\caption{Incumbency outcomes: \textsc{bisect} vs.\ enacted 2022 maps vs.\ random
baseline. All \textsc{bisect} results are single-run (seed\,$=42$).}
\label{tab:incumbency-summary}
\begin{tabular}{l l r r r}
\toprule
State & Metric & \textsc{bisect}$^\dagger$ & Enacted 2022 & Random Baseline \\
\midrule
NC ($k=14$) & Pairing rate & 0.141 & 0.000 & 0.222 \\
            & Safe-seat fraction & 0.43  & 0.64  & n/a \\
            & Open-seat count & 4     & 2     & 1.0--2.5 \\
\addlinespace
WI ($k=8$)  & Pairing rate & 0.107 & 0.000 & 0.175 \\
            & Safe-seat fraction & 0.38  & 0.50  & n/a \\
            & Open-seat count & 2     & 1     & 0.8--1.6 \\
\addlinespace
TX ($k=38$) & Pairing rate & 0.105 & 0.026 & 0.153 \\
            & Safe-seat fraction & 0.55  & 0.71  & n/a \\
            & Open-seat count & 7     & 3     & 4.6--6.1 \\
\bottomrule
\end{tabular}
\begin{tablenotes}
\small
\item[$\dagger$] Single-run result (seed\,$=42$). See I.1--I.3 for ensemble distributions.
\item The random baseline for safe-seat fraction is not reported here because it depends
  on geographic partisan sorting, not on a simple combinatorial formula; see I.2 for the
  geographic-constraint-adjusted baseline.
\end{tablenotes}
\end{threeparttable}
\end{table}

\subsection{Incumbent Pairing}

The most striking finding in Table~\ref{tab:incumbency-summary} is the enacted maps'
pairing rates. North Carolina and Wisconsin enacted maps achieve a pairing rate of
exactly 0.000---every sitting incumbent received a home district under the enacted plan.
Texas, with 38 districts and 36 incumbents seeking reelection, achieves a pairing rate
of 0.026 (approximately one pairing event), substantially below the random baseline of
0.153. The \textsc{bisect} algorithmic maps produce pairing rates of 0.141 (NC),
0.107 (WI), and 0.105 (TX), all materially closer to the random baseline than to the
enacted map rates.

As I.1 documents in detail, this pattern is consistent with the hypothesis that enacted
maps use incumbent address data to draw boundaries that guarantee each sitting member
a home district, while \textsc{bisect} maps produce pairing at the rate that geographic
chance implies.

\subsection{Safe-Seat Fraction}

Across all three states, \textsc{bisect} maps produce fewer safe seats than enacted maps.
The gap is largest in North Carolina (0.43 vs.\ 0.64) and Texas (0.55 vs.\ 0.71) and
more modest but still present in Wisconsin (0.38 vs.\ 0.50). These differences are directionally consistent across all three states; I.2 reports
$p < 0.01$ for NC and TX under a two-sample test of proportions using single-seed bisect
fractions ($^\dagger$), which are indicative but require multi-seed ensemble validation
before interpretation as a general property of the algorithm.

The safe-seat reduction is mechanically driven by \textsc{bisect}'s compactness
optimisation: compact geographic districts in states with moderate partisan segregation
are less likely to produce extreme partisan margins than districts drawn with an eye to
concentrating one party's voters in safe enclaves and dispersing the other party's
voters across competitive seats.

\subsection{Open-Seat Count}

\textsc{bisect} maps produce more open seats than enacted maps in all three states.
NC produces 4 bisect open seats versus 2 enacted; WI produces 2 versus 1; TX produces
7 versus 3. The \textsc{bisect} open-seat counts are within or slightly above the
random-baseline ranges, consistent with the view that algorithmic indifference to
incumbency produces open seats at approximately the rate that geographic chance
implies---not more disruptive than a random process, but substantially more open than
enacted maps that actively engineer home districts for all sitting members.
